\subsection{Uniform perturbation in primitive quotients}
\label{sec:perturbation}

We now turn the uniform gap into a perturbation that works in many
primitive quotients at once.  The fixed amplification condition
\propF{} supplies finite reducing copies on a spectral cut; the gap
then prevents one fixed element from nearly cancelling every such
copy.  We first record the two consequences of the comparison and
transport results in Section~\ref{sec:foundations} that this argument
uses.

\begin{lemma}[Fixed-amplification ideal comparison]
\label{lem:ideal-comparison}
Let \(C\) be a \(C^*\)-algebra, let \(h\in C_+\), and let \(n\geq1\)
be an integer such that \(h^{\oplus n}\) is properly infinite.
For every \(x\in\Id(h)_+\) and every integer \(k\geq1\),
\begin{equation}
 x^{\oplus k}\precsim h^{\oplus n}.
 \label{eq:ideal-multiplicity}
\end{equation}
In particular, if \(C\) satisfies \propF{}, this conclusion holds for
every \(h\in C_+\) with the same \(n\).
\end{lemma}

\begin{proof}
In the common stabilization,
\[
 \Id_{C\otimes\K}(h^{\oplus n})
   =\Id_C(h)\otimes\K.
\]
Thus \(x^{\oplus k}\) belongs to the ideal generated by
\(h^{\oplus n}\).  The proper-target ideal comparison in
Lemma~\ref{lem:comparison-rules} gives
\eqref{eq:ideal-multiplicity}.  This argument takes place in the
stabilization and therefore covers every finite rectangular corner
and nonunital \(C\).  It also covers \(h=0\), when \(x=0\).
\end{proof}

\begin{lemma}[Exact implementation of a spectral cut]
\label{lem:cut-implementation}
Let \(C\) be a \(C^*\)-algebra, let \(m,s\geq1\), and suppose that
\(x\in M_m(C)_+\) and \(h\in M_s(C)_+\) satisfy \(x\precsim h\).
If \(f\colon[0,\|x\|]\to[0,\infty)\) is continuous and vanishes
on a neighborhood of zero, then there is \(R\in M_{s,m}(C)\) such that
\begin{equation}
 R^*R=f(x),\qquad RR^*\in\Her(h),\qquad
 \|R\|^2=\|f(x)\|.
 \label{eq:exact-cut}
\end{equation}
\end{lemma}

\begin{proof}
If \(f(x)=0\), take \(R=0\).  Otherwise, choose \(\delta>0\) so
that \(f\) vanishes on \([0,\delta]\), and choose
\(0<\varepsilon<\delta\).  Then
\(f(x)^{1/2}\in\Her((x-\varepsilon)_+)\).
Apply Lemma~\ref{lem:transport} to \(x\precsim h\) at this cutoff,
and let \(w\in M_{s,m}(C^{**})\) be its rectangular partial isometry.
Its initial projection is the support of \((x-\varepsilon)_+\).
The membership and support conclusions of that lemma therefore give
\[
 R=w f(x)^{1/2}\in M_{s,m}(C),\qquad
 R^*R=f(x),\qquad RR^*=w f(x)w^*\in\Her(h).
\]
The norm identity follows from \(\|R\|^2=\|R^*R\|\).
Only the product \(R\), rather than \(w\) itself, is asserted to
belong to the rectangular algebra.
\end{proof}

For \(P\in\Prim(C)\), a subscript \(P\) denotes the image in \(C/P\),
also at matrix levels.  An irreducible representation with kernel
\(P\) faithfully represents this quotient and computes its norm.
The primitive ideal test in Lemma~\ref{lem:ideals-cutoffs} will let us
recognize ideal membership from these quotient images.

The uniform gap will be applied to the unitization of a matrix algebra
over \(C\).  We record the two permanence facts that make this a valid
exact target, also when \(C\) has no unit.

\begin{lemma}[The exact target for matrix models]
\label{lem:exact-matrix-unitization}
If \(A\) is an exact \(C^*\)-algebra and \(n\ge1\), then
\(\widetilde{M_n(A)}\) is exact.
\end{lemma}

\begin{proof}
Finite matrix exactness follows, for every \(C^*\)-algebra \(D_0\),
from the identifications
\[
 M_n(A)\otimes_{\min}D_0=M_n(A\otimes_{\min}D_0),
\]
which preserve the ideal kernels in the definition of exactness.
For the unitization step, let \(A_0\) be a nonunital exact algebra,
let \(I\triangleleft D_0\) be a closed two-sided ideal, and write
\(q_0\colon D_0\to D_0/I\).  The scalar character
\(\chi\colon\widetilde A_0\to\C\) gives, for every
\(z\in\widetilde A_0\otimes_{\min}D_0\), the bounded decomposition
\[
 z=z_0+1\otimes d,\qquad
 d=(\chi\otimes\id)(z),\quad z_0\in A_0\otimes_{\min}D_0.
\]
Indeed, both projections are bounded, and the decomposition holds
first on algebraic tensors and then on their completion.  If
\((\id\otimes q_0)(z)=0\), applying \(\chi\otimes\id\) gives
\(d\in I\).  The isometric inclusion
\[
 A_0\otimes_{\min}(D_0/I)
 \subset\widetilde A_0\otimes_{\min}(D_0/I)
\]
then shows \((\id\otimes q_0)(z_0)=0\).  Exactness of \(A_0\) gives
\(z_0\in A_0\otimes_{\min}I\), hence
\(z\in\widetilde A_0\otimes_{\min}I\).  If \(M_n(A)\) is already
unital, this second step is unnecessary.
\end{proof}

Fix, once and for all, a constant
\begin{equation}
 0<\sigma<\min\{1,\gamma/3\},
 \label{eq:perturbation-constant}
\end{equation}
where \(\gamma\) is supplied by Proposition~\ref{prop:uniform-gap}.
This \(\sigma\) is independent of the algebra and of the fixed
amplification integer.

\begin{proposition}[Uniform perturbation]\label{prop:perturbation}
Let \(C\) be an exact \(C^*\)-algebra satisfying \propF{} for one
fixed integer \(n\geq1\).  Let \(p,q\in C\) be contractions and
let \(h,k,c\in C_+\) satisfy
\begin{equation}
 p^*ph=h,\qquad qq^*k=k,\qquad
 c^{\oplus n}\precsim h,\qquad c^{\oplus n}\precsim k.
 \label{eq:perturbation-hypotheses}
\end{equation}
For every \(T\in C\) and every \(\eta,\alpha>0\), there is a
contraction \(y\in C\) such that
\begin{equation}
 \|(T+\eta pyq)_P\|\geq\sigma\eta
 \quad\text{for every }P\in\Prim(C)\text{ with }\|c_P\|\geq\alpha.
 \label{eq:perturbation-conclusion}
\end{equation}
If \(h,k\) belong to a hereditary subalgebra
\(\mathcal D\subset C\), then \(y\) may be chosen in \(\mathcal D\).
If, in addition or independently, \(p,q\) are positive and \(pq=0\),
then \(y\) may be chosen selfadjoint.  Both conclusions may be
imposed simultaneously.  The constant \(\sigma\) is independent
of all the data in the Proposition, including \(n\) and the norms
of \(T,h,k,c\).
\end{proposition}

\begin{proof}
Write \(K=\{P\in\Prim(C):\|c_P\|\geq\alpha\}\).  If \(K\) is
empty, \(y=0\) satisfies all conclusions.  Suppose henceforth that
\(K\) is nonempty.  The contradiction will have three levels of choices.
First we fix comparison rows and a single contraction in \(M_n(C)\).  Each finite matrix contraction then determines an
admissible perturbation.  If it fails, a primitive quotient witnessing
that failure supplies an approximate reducing copy for that fixed
contraction.  All these copies therefore come from one element
of the same exact algebra \(\widetilde{M_n(C)}\).

\smallskip
\noindent\emph{Fixed comparison rows.}
Choose a continuous function \(e\colon[0,\infty)\to[0,1]\) that
vanishes near zero and is one on \([\alpha/2,\infty)\).
Apply Lemma~\ref{lem:cut-implementation} to the last two comparisons
in \eqref{eq:perturbation-hypotheses}.  This gives rows
\(R,S\in M_{1,n}(C)\) with
\begin{equation}
 R^*R=S^*S=E=e(c)^{\oplus n},\qquad
 RR^*\in\Her(h),\quad SS^*\in\Her(k),\qquad
 \|R\|,\|S\|\leq1.
 \label{eq:perturbation-rows}
\end{equation}
The relation \(p^*ph=h\), together with its adjoint, shows that
left and right multiplication by \(p^*p\) are the identity on
\(\Her(h)\).  Hence
\((1-p^*p)RR^*(1-p^*p)=0\), which implies
\((1-p^*p)R=0\).  The same argument for \(S\) gives
\begin{equation}
 p^*pR=R,\qquad qq^*S=S.
 \label{eq:perturbation-row-identities}
\end{equation}
At this point fix the single element
\begin{equation}
 V=\eta^{-1}R^*p^*Tq^*S\in M_n(C).
 \label{eq:fixed-testing-element}
\end{equation}
The rows \(R,S\) and the element \(V\) will not depend on any finite
matrix model.  Fix now also the global contraction
\[
 \widehat V=V(\max\{1,|V|\})^{-1}\in M_n(C).
\]
The inverse is evaluated in the unitization, so its product with
\(V\) remains in \(M_n(C)\).

\smallskip
\noindent\emph{Reducing copies of arbitrary finite contractions.}
Let \(d\geq1\) and let \(D\in M_d(\C)\) be a contraction.
Put \(c_1=(c-\alpha/2)_+\).  Condition \propF{} and
Lemma~\ref{lem:ideal-comparison}, applied to the scalar element
\(c_1\), give
\[
 c_1^{\oplus d}\precsim c_1^{\oplus n}.
\]
Choose a continuous \(\varphi\colon[0,\infty)\to[0,1]\), vanishing
near zero and equal to one on \([\alpha/4,\infty)\), and put
\(e_1(c)=\varphi(c_1)\).  Lemma~\ref{lem:cut-implementation}
yields \(L\in M_{n,d}(C)\) such that
\begin{equation}
 L^*L=e_1(c)^{\oplus d},\qquad
 LL^*\in\Her(c_1^{\oplus n}),\qquad \|L\|\leq1.
 \label{eq:finite-model-column}
\end{equation}
Because \(e(c)c_1=c_1\), multiplication by \(E\) is the identity
on \(\Her(c_1^{\oplus n})\).  The range-square argument used for
\eqref{eq:perturbation-row-identities} gives \(EL=L\).
Define
\begin{equation}
 w=LDL^*\in M_n(C),\qquad \|w\|\leq1,\qquad EwE=w.
 \label{eq:finite-model}
\end{equation}
Scalar matrix multiplication here requires no unit in \(C\).

For each \(P\in K\), choose an irreducible representation
\(\pi_P\colon C\to B(H_P)\) with kernel \(P\), and use its matrix
amplifications.  Identify \(C/P\) with the image of the induced
faithful representation on \(H_P\), and make the same identification
at rectangular matrix levels.  The spectral subspace
\(1_{(3\alpha/4,\infty)}(\pi_P(c))H_P\) is nonzero, since
\(\|c_P\|\geq\alpha\).  Choose a unit vector \(\xi_P\) in that
subspace, so that \(\pi_P(e_1(c))\xi_P=\xi_P\).  Define
\[
 J_P\colon\C^d\longrightarrow H_P^{\oplus n},\qquad
 J_P\lambda=L_P(\lambda\otimes\xi_P).
\]
Equation~\eqref{eq:finite-model-column} proves \(J_P^*J_P=1_d\).
Both \(D\) and \(D^*\) preserve \(\C^d\otimes\xi_P\), and the
same equation gives
\begin{equation}
 w_PJ_P=J_PD,\qquad w_P^*J_P=J_PD^*.
 \label{eq:reducing-model}
\end{equation}
Thus one copy reduces both \(w_P\) and its adjoint.  No eigenvector
at the norm of \(c_P\), and no continuous choice of the vectors
\(\xi_P\), is required.

\smallskip
\noindent\emph{Admissible tests and the two optional conclusions.}
For the chosen \(d,D\), take
\[
 y_0=-RwS^*\in C,\qquad\|y_0\|\leq1.
\]
If \(h,k\in\mathcal D\), then \(RR^*,SS^*\in\mathcal D\).
For a positive contractive approximate unit \((a_\lambda)\) of
\(\mathcal D\), the identity
\[
 \|(1-a_\lambda)R\|^2
 =\|(1-a_\lambda)RR^*(1-a_\lambda)\|\longrightarrow0
\]
and its analogue for \(S\) give
\(a_\lambda RwS^*a_\lambda\to RwS^*\) in norm.
Each approximant belongs to
\(\mathcal D C\mathcal D\subset\mathcal D\), so \(y_0\in\mathcal D\).
This uses the range squares of the rows and makes no assertion that
their individual entries belong to \(\mathcal D\).

If \(p,q\geq0\) and \(pq=0\), then \(p^2R=R\) implies \(pR=R\):
multiply \((1-p^2)R=0\) by \((1+p)^{-1}\) in the unitization.
Likewise \(qS=S\).  It follows that \(pS=qR=0\) and \(S^*R=0\).
Consequently \(y_0^2=0\) and \(py_0^*q=0\).  The positive elements
\(y_0y_0^*\) and \(y_0^*y_0\) are orthogonal, whence
\[
 \|y_0+y_0^*\|^2
 =\|y_0y_0^*+y_0^*y_0\|=\|y_0\|^2\leq1.
\]
We may therefore use \(y=y_0+y_0^*\) when selfadjointness is
requested, and \(y=y_0\) otherwise.  In both cases \(pyq=py_0q\);
hereditary membership is preserved if it is also requested.

\smallskip
\noindent\emph{One global contraction and the uniform gap.}
Suppose that no admissible \(y\) satisfies
\eqref{eq:perturbation-conclusion}.  For each \(d,D\), choose a
witness \(P\in K\) for the test just constructed.  Then
\[
 \|(T-\eta pRwS^*q)_P\|<\sigma\eta.
\]
Multiplying by the contractive outer factors \(R^*p^*\) and \(q^*S\),
and using \eqref{eq:perturbation-row-identities}, gives
\begin{equation}
 \|(V-EwE)_P\|=\|(V-w)_P\|<\sigma.
 \label{eq:local-near-model}
\end{equation}
Although \(P\) may depend on \(d,D\), the element \(V\) is the one
fixed in \eqref{eq:fixed-testing-element}.

Functional calculus commutes with quotient maps for the fixed global
clipping \(\widehat V\).  Hence in each witness quotient
\begin{equation}
 \|V_P-\widehat V_P\|
 =\|(|V_P|-1)_+\|
 =(\|V_P\|-1)_+<\sigma.
 \label{eq:quotient-clipping}
\end{equation}
The last strict inequality follows from \eqref{eq:local-near-model}
and \(\|w_P\|\leq1\).  The clipping identity follows by squaring the
absolute value of
\(V_P-V_P\max\{1,|V_P|\}^{-1}\) and applying scalar functional
calculus.  Combining \eqref{eq:local-near-model},
\eqref{eq:quotient-clipping}, and \eqref{eq:reducing-model} gives
\begin{equation}
 \|\widehat V_PJ_P-J_PD\|<2\sigma<\gamma,\qquad
 \|\widehat V_P^*J_P-J_PD^*\|<2\sigma<\gamma.
 \label{eq:forbidden-intertwining}
\end{equation}

By Lemma~\ref{lem:exact-matrix-unitization},
\(B=\widetilde{M_n(C)}\) is a unital exact algebra.
The amplified representations \(\pi_P^{(n)}\) extend unitally to
\(B\).  Although the witnessing \(P\) varies with \(d,D\), the
contraction \(\widehat V\) and the exact algebra \(B\) remain fixed.
Thus \(\widehat V\in B\) has, for every \(d,D\), a unital
representation and an isometry satisfying
both inequalities in \eqref{eq:forbidden-intertwining}.  This
contradicts Proposition~\ref{prop:uniform-gap} and proves the
Proposition with both optional conclusions and the same \(\sigma\).
\end{proof}
